% filename: chap-introduction-v2.tex
% Strategy: §1.1 Background and §1.2 Current Practices preserved from Phase-I draft
%   (rich industry stats and citations intact).
%   §1.3 onward rewritten to:
%     - remove all Phase-I / Phase-II framing
%     - replace "perishable inventory" with "finite replenishable inventory"
%     - replace SC-PPO with MaskablePPO / constrained RL
%     - correct single-site scope claim to multi-site evaluation
%     - correct 15-minute → hourly (Δt = 1h, consistent with Ch3/Ch5)
%     - update thesis organisation to match actual chapter content
%     - add H1/H2/H3 hypotheses and key empirical numbers
% ============================================================

\chapter{Introduction}
\label{chap:introduction}

% ============================================================
\section{Background and Motivation}
% ============================================================

Information and Communication Technology (ICT) has become a foundational
pillar of the global digital economy, powering mobile connectivity, cloud
services, fintech, digital governance, and critical public infrastructure.
As of 2025, mobile networks serve more than \textbf{5.6--5.8 billion unique
subscribers} worldwide, representing over 70\% of the global population
\citep{gsma2025_mobile_economy}. Mobile data traffic continues to grow at
double-digit annual rates, driven by video streaming, Internet-of-Things
(IoT) proliferation, and the ongoing densification of 4G and rollout of 5G
networks, including in rural and semi-urban regions of emerging markets.

At the core of this connectivity lies the \textbf{telecom tower site} (or base
transceiver station, BTS site), which must operate 24\,$\times$\,7 to ensure
uninterrupted service. A typical multi-technology tower consumes
\textbf{2--8\,kW continuously}, making the Radio Access Network (RAN)
responsible for \textbf{60--80\%} of a mobile operator's total electricity
demand \citep{americas2023_energy_efficiency}. Globally, mobile networks
already account for approximately \textbf{2--3\%} of worldwide electricity
consumption, a figure projected to rise further as 5G sites increase per-site
power demand by \textbf{2--3$\times$} compared to 4G
\citep{mckinsey2023_greener_telecom}.

In developing regions like India and Sub-Saharan Africa, telecom towers
commonly face unreliable grid access and rely heavily on diesel backup
\citep{gsma2020_tower_esco}. To guarantee high network uptime, operators
rely on \textbf{diesel generators (DGs)} and \textbf{battery banks}.
\textbf{Diesel fuel represents 25--30\% of tower OPEX, rising to 70\% at
remote sites} \citep{datoms2024,gsma2020_tower_esco}.
In India, 60--70\% of towers rely on diesel due to inconsistent grid supply
\citep{datoms2024}.

\textbf{Fuel logistics failures---through pilferage, stockouts, and delivery
disruptions---pose significant risks to tower reliability}; the 2025 threat
to 16{,}000 Nigerian sites underscores diesel supply-chain vulnerability
\citep{datoms2024,alton2025}. The heavy reliance on diesel generators in
weak-grid regions therefore drives both high operational costs and substantial
carbon emissions. In the Indian context, a typical macro BTS site that relies
on a diesel generator for roughly eight hours per day consumes about
8{,}760\,L of diesel per year, corresponding to approximately
\textbf{23--24\,t CO\textsubscript{2} per site per year} when using a standard
emission factor of 2.68\,kg\,CO\textsubscript{2} per litre of diesel
\citep{trai2011_green_telecom,ghgprotocol2017_emission_factors}.

These operational and environmental challenges are intensified by global
sustainability commitments. The GSMA Climate Action roadmaps and major
operators target net-zero emissions by 2040--2050, yet diesel-powered sites
remain one of the largest controllable sources of Scope\,1 emissions
\citep{gsma2023}. Hybrid renewable systems (solar PV + battery + diesel) have
demonstrated \textbf{20--65\% reductions} in diesel consumption and CO\textsubscript{2}
emissions in field deployments \citep{gsma2020_tower_esco,datoms2024}.
However, maximising these savings amid stochastic solar irradiance,
traffic-driven loads, and grid outages requires sophisticated control beyond
current rule-based systems.

In India---the world's second-largest mobile market with over
\textbf{1.2 billion subscribers} and approximately
\textbf{400{,}000--500{,}000 tower sites}---the challenge is particularly
acute. Rural and semi-urban grid availability often falls below 14 hours per
day, making diesel dominant at many sites \citep{trai2024,datoms2024}.
Aggressive 5G rollout and renewable energy mandates intensify pressure for
low-carbon hybrid transitions while controlling OPEX.


% ============================================================
\section{Current Practices and Their Limitations}
% ============================================================

Industry-standard energy management today relies predominantly on
\textbf{rule-based} or \textbf{threshold-driven controllers} embedded in
Power Management Units (PMUs) or Hybrid Power Controllers (HPCs). Typical
rules include: use solar whenever available; switch to battery when the grid
fails; and start the diesel generator if battery State-of-Charge (SoC) falls
below 30\%.

These heuristics are simple, robust, and inexpensive, but they are
fundamentally \emph{myopic}: they ignore long-term fuel cost, battery
degradation cost, forecasted outages, variable diesel prices, and potential
carbon penalties. Tools such as \textbf{HOMER Pro} are widely used for system
design and planning. They can produce static daily dispatch schedules and size
PV/DG/battery systems that reduce diesel use by 30--58\%, yet they operate
offline, typically under deterministic or average conditions, and cannot adapt
in real time to uncertainty \citep{homer_energy_2020,meena2020_homer_telecom}.

More advanced approaches---such as Model Predictive Control (MPC), stochastic
programming, and classical optimisation---improve performance by explicitly
optimising over a rolling horizon. However, they demand accurate short-term
forecasts and non-trivial on-site computation, making them less attractive for
low-cost tower controllers in remote locations.

Reinforcement Learning (RL) has recently shown strong results in microgrids,
electric vehicles, and building energy management, where agents learn to trade
off cost and reliability under uncertainty
\citep{mao2021_inventory_rl,lin_ev_rl_2024,wang2022_building_rl}. Recent
work specifically on telecom base station energy management has demonstrated
18--47\% reductions in operating cost relative to rule-based and MPC baselines
\citep{ma_pan_2025_telecom_ies,yan_2025_safe_marl_bs,alhajhassan_2024_battery_ageing_rl}.
Yet a decisive modelling gap persists across this body of work: diesel fuel is
universally treated as an always-available, unconstrained resource. To the best
of the author's knowledge, no prior RL formulation for telecom towers models
the diesel tank level as a finite depletable stock or treats replenishment
ordering as a control action. Existing formulations also do not account for
stochastic delivery lead times and their interaction with dispatch decisions.
This gap motivates the
methodology adopted in this thesis. The following chapter reviews related
work across telecom energy management, microgrid scheduling, and
inventory-aware control, and positions the proposed formulation relative
to existing approaches.


% ============================================================
\section{Research Problem and Objectives}
% ============================================================

This thesis addresses the problem of how to optimally manage energy resources
and diesel replenishment in hybrid telecom towers under significant uncertainty
in solar generation, grid availability, and diesel delivery, so as to minimise
unserved energy and operational expenditure while ensuring service continuity.

The central research question is:

\begin{quote}
\emph{Does jointly learning diesel ordering and energy dispatch within a single
reinforcement-learning agent reduce unserved energy at remote telecom
base stations relative to approaches that treat these decisions separately?}
\end{quote}

To answer this question, the thesis formulates the problem as a
\textbf{Constrained Markov Decision Process (CMDP)} and evaluates three
hypotheses. The empirical evaluation focuses on reliability
measured through \textbf{Expected Energy Not Served (EENS)} and related
operational metrics:

\begin{description}
  \item[H1 (Joint ordering).] Does joint learning of dispatch and diesel
  ordering reduce unserved energy relative to a decoupled $(s,S)$ ordering
  policy with separately learned dispatch?
  \item[H2 (Action masking).] Does hard action masking reduce constraint
  violations and improve training stability relative to unconstrained RL?
  \item[H3 (Inventory observation).] Does including the diesel inventory state
  in the agent's observation space improve reliability relative to a
  dispatch-only policy without fuel-level awareness?
\end{description}

The specific objectives of this work are to:

\begin{enumerate}
  \item Formulate the hybrid telecom tower energy management problem as a CMDP
  that treats battery state-of-charge and diesel fuel level as coupled
  state variables under uncertain supply and demand, with diesel
  modelled as a \textbf{finite replenishable inventory} subject to stochastic
  delivery lead times.

  \item Design and train reinforcement learning agents capable of jointly
  learning energy dispatch and diesel replenishment strategies under
  operational constraints, and a hierarchical decomposition baseline in which
  ordering and dispatch are learned separately.

  \item Develop a Gymnasium-compatible simulation environment driven by real
  ITU/Zindi telecom site traces that reproduces empirically grounded load,
  PV generation, and grid availability dynamics at hourly resolution.

  \item Benchmark the learned RL policies against rule-based and
  decomposition-based baselines through a structured eight-variant ablation
  study, quantifying improvements in unserved energy, diesel consumption, and
  network uptime across multiple operating scenarios.
\end{enumerate}


% ============================================================
\section{Contributions}
% ============================================================

The main contributions of this thesis are:

\begin{enumerate}
  \item A telecom-tower-specific, inventory-theoretic CMDP formulation that
  explicitly represents battery state-of-charge and diesel fuel volume as
  state variables subject to uncertain demand, stochastic replenishment lead
  times, and cost/emission penalties.

  \item The design and empirical evaluation of a constrained reinforcement
  learning framework---\textbf{RLInv}---for hybrid telecom tower energy
  management, including inventory-aware state and action design, with hard
  feasibility enforcement via action masking, and a hierarchical decomposition
  baseline (TrackB) using a classical $(s,S)$ ordering policy.

  \item An open, reproducible simulation environment and experimental protocol
  integrating the ITU~2024 Smart Energy Scheduling dataset (ten base station
  sites, 60~days each) with a synthetic lead-time overlay, supporting
  reproducible multi-site evaluation.

  \item An empirical evaluation across two logistics scenarios and eight policy
  variants, demonstrating substantial reductions in unserved energy relative to
  the hierarchical baseline---including a 49.6\% improvement under the trained
  logistics scenario---and characterising the robustness limits that motivate
  multi-scenario training in future work.
\end{enumerate}


% ============================================================
\section{Scope and Limitations}
% ============================================================

To maintain focus and analytical tractability, this thesis makes the following
scope choices:

\begin{itemize}
  \item The control problem is formulated at the level of an individual tower
  site; empirical evaluation is conducted across all ten ITU/Zindi sites to
  assess site-to-site robustness and heterogeneity.
  \item Control decisions are taken at \textbf{hourly resolution}
  ($\Delta t = 1$\,h), consistent with the temporal granularity of the
  ITU/Zindi dataset.
  \item Component sizes (PV, battery, DG rating, tank capacity) are assumed
  given; the focus is on \emph{operational control}, not investment planning.
  \item Standard empirical models are used for battery ageing and generator
  efficiency; detailed electrochemical degradation models are outside the scope.
  \item The in-transit diesel state variable ($\mathrm{Pipe}_t$) is the sole
  synthetic component, constructed via a \textbf{synthetic stochastic lead-time
  overlay applied to the diesel delivery process}; all other inputs are drawn
  directly from the ITU site traces.
  \item Statistical comparisons use $n = 10$ independent seeds with paired
  testing across identical site--scenario combinations (an increase from
  an $n=3$ pilot evaluation conducted earlier in the project). This
  improves statistical power considerably, but some comparisons ---
  particularly the inventory-observability ablation (H3) --- remain
  underpowered to detect small effects with high confidence; further
  replication would strengthen these findings.
\end{itemize}

Limitations arising from these choices include:

\begin{itemize}
  \item The RL agent is trained under a single logistics scenario and does not
  generalise robustly to delayed delivery conditions.
  \item Multi-site coordination and shared fuel logistics are not modelled.
  \item Integrated forecasting modules for load or solar are not learned jointly
  with the control policy.
  \item The diesel generator is modelled at fixed rated power when ON; continuous
  power setpoint control is left for future work.
\end{itemize}

These aspects are identified as promising directions for future work.
% ============================================================
\section{Thesis Roadmap and Relation Between Chapters}
% ============================================================

To make the logical flow of the thesis explicit,
Fig.~\ref{fig:thesis-roadmap} summarises how the chapters relate to
one another and highlights where the primary contributions arise. The
thesis progresses from motivation and gap identification, through formal
problem definition and implementation, to empirical evaluation and
conclusions. The novelty of the work lies in connecting three elements
that are typically treated separately in prior literature:
(i)~telecom tower hybrid-energy dispatch,
(ii)~diesel inventory and replenishment uncertainty, and
(iii)~constrained reinforcement learning for joint operational control.
Chapters~3, 4, and~5 form the core technical contributions of the thesis,
while Chapter~7 integrates these components into a unified empirical
evaluation.

\begin{figure}[!ht]
\centering
\resizebox{0.82\textwidth}{!}{%
\begin{tikzpicture}[
    node distance=0.4cm and 1.6cm,
    every node/.style={font=\small, align=center},
    box/.style={
        draw, rounded corners, rectangle,
        minimum width=3.4cm, minimum height=1.05cm
    },
    novel/.style={
        draw, rounded corners, rectangle,
        minimum width=3.4cm, minimum height=1.05cm,
        fill=gray!15, thick
    },
    arrow/.style={->, thick},
    xarrow/.style={->, thick, dashed}
]
\node[box]   (c1) {Chapter 1\\Motivation, scope,\\research objectives};
\node[box,   below=of c1] (c2) {Chapter 2\\Literature review\\and research gap};
\node[novel, below=of c2] (c3) {\textbf{Chapter 3}\\CMDP: joint dispatch\\+ inventory formulation};
\node[novel, below=of c3] (c4) {\textbf{Chapter 4}\\Methodology: TelecomEnv,\\RLInv, MaskablePPO};
\node[novel, below=of c4] (c5) {\textbf{Chapter 5}\\Dataset, EDA, site\\difficulty classification};
\node[box,   below=of c5] (c6) {Chapter 6\\Experimental execution\\and training runs};
\node[box,   below=of c6] (c7) {Chapter 7\\Results, ablations,\\robustness analysis};
\node[box,   below=of c7] (c8) {Chapter 8\\Conclusion, limitations,\\future work};

\draw[arrow] (c1) -- (c2);
\draw[arrow] (c2) -- (c3);
\draw[arrow] (c3) -- (c4);
\draw[arrow] (c4) -- (c5);
\draw[arrow] (c5) -- (c6);
\draw[arrow] (c6) -- (c7);
\draw[arrow] (c7) -- (c8);

\draw[xarrow] (c3.east) to[out=0, in=0, looseness=1.4]
    node[right, font=\footnotesize] {CMDP defines\\eval metrics}
    (c7.east);
\draw[xarrow] (c5.east) to[out=0, in=0, looseness=0.9]
    node[right, font=\footnotesize] {ITU data\\grounds results}
    (c7.east);

\node[draw, dashed, rounded corners, fit=(c3)(c4)(c5),
      inner sep=4pt,
      label={[font=\footnotesize\itshape]right:Novel\\contributions}] {};
\end{tikzpicture}%
}
\caption{Thesis roadmap. Shaded boxes (Chapters~3--5) represent the
primary novel contributions. Dashed arrows indicate cross-chapter
dependencies. The vertical spine shows the sequential reading order.}
\label{fig:thesis-roadmap}
\end{figure}
\FloatBarrier
% ============================================================
\section{Thesis Organisation}
% ============================================================

Following the roadmap shown in Fig.~\ref{fig:thesis-roadmap}, the
remainder of this thesis elaborates each stage of the research workflow
in detail:

\begin{itemize}
  \item \textbf{Chapter~\ref{chap:literature}} reviews energy management
  practices in telecom towers, hybrid PV--battery--DG systems, inventory
  theory in fuel and energy logistics, and RL applications in energy systems.
  It identifies the gap that motivates this work and positions the proposed
  framework relative to prior art.

  \item \textbf{Chapter~\ref{chap:problem}} defines the hybrid telecom tower
  system model, the inventory-theoretic CMDP formulation, the state and action
  spaces, the cost function, and the hard operational constraints.

  \item \textbf{Chapter~\ref{chap:methodology}} presents the implemented
  solution framework: the \texttt{TelecomEnv} simulator, the \texttt{MaskablePPO}
  agent architecture, the $(s,S)$~hierarchical baseline, the eight ablation
  variants, the training pipeline, and the evaluation protocol.

  \item \textbf{Chapter~\ref{chap:datasets-setup}} documents the ITU/Zindi
  dataset, preprocessing steps, cross-site heterogeneity, and the exploratory
  data analysis that informed simulator design and site selection for the
  main experimental evaluation.

  \item \textbf{Chapter~\ref{chap:experimental-execution}} records the
  execution of the experimental plan: training runs, convergence monitoring,
  and evaluation campaign details.

  \item \textbf{Chapter~\ref{chap:results}} presents and interprets the full
  experimental results, including the main EENS comparison, the structured
  ablation addressing H1, H2, and~H3, the delayed-logistics robustness
  analysis, and cross-site generalisation findings.

  \item \textbf{Chapter~\ref{chap:conclusion}} summarises the principal
  contributions and limitations, and outlines directions for future work
  including multi-scenario training and expanded statistical evaluation.
\end{itemize}

\clearpage